\section{投资问题：H1高利润有多少来自可持续能力}

文莱一期是本轮估值重估的核心。投资问题不是装置是否满产，而是高负荷、区域供给缺口、产品裂解价差、原油贴水、运费、税惠与70\%权益结合后，能够形成多少可持续归母现金利润。核心判断是，2026H1利润跃升具有经营基础，但峰值价差和税后归属尚未得到分部财务拆分，不能把H1运行率外推到全年和终值。

文莱一期设计原油加工能力800万吨／年，主要产品包括565万吨／年成品油和PX／苯合计265万吨／年。公司称原油主要来自文莱和马来西亚，采购价格以Brent为结算基准，产品主要销往东南亚和澳大利亚。相较中国沿海大炼化，文莱的差异不是绝对规模，而是东南亚区域供需、较短采购半径和先锋企业税惠。

\begin{exhibitbox}[文莱一期资产与经济变量]
\noindent\begin{tabularx}{\textwidth}{L{3.15cm}L{2.4cm}L{2.8cm}X}
\toprule
\textbf{项目} & \textbf{已披露} & \textbf{盈利变量} & \textbf{未披露边界} \\
\midrule
原油加工 & 800万吨／年 & 进料量、收率、检修、能耗 & 2026实际进料与综合收率 \\
成品油 & 565万吨／年 & 柴油／航煤裂解、区域供需 & 分产品销量、ASP与客户 \\
芳烃 & PX／苯合计265万吨／年 & PX--石脑油、苯价差 & PX／苯拆分量价成本 \\
原油来源 & 文莱、马来西亚等；Brent基准 & 贴水、品质、库存与运输 & 供应商、长约／现货、实际贴水与运费 \\
销售区域 & 东南亚、澳大利亚 & 区域缺口与出口物流 & 具名客户、地区销量与信用期 \\
权益／税 & 上市公司70\%；先锋企业11年免税 & 少数股东、有效税率和到期日 & 税惠具体起止、2026少数股东拆分 \\
\bottomrule
\end{tabularx}
\end{exhibitbox}
\sourcenote{HY-OFF-2025AR；IND-BRUNEI；财务与项目资料截止2026-07-22。公司披露的区域与产品方向不等于具名客户或2026实际销量。}

这张表说明文莱的利润并非由单一柴油价差决定。即使产品裂解走强，若原油贴水收窄、运费上升、装置检修或美元融资成本增加，综合净利也可能低于行业代理暗示；反之，较低原油成本和税惠可使项目利润高于单一裂解价差的表面水平。

\section{2025低基数：项目已修复，但尚未展现峰值能力}

恒逸文莱单体2025年收入404.74亿元、净利润2.15亿元；分析包以合并主体口径归纳收入421.40亿元、净利润2.42亿元、经营现金流3.25亿元。不同表格可能对应单体与主体汇总口径，本报告不将两组数混用：正文以年报主要子公司单体表呈现404.74／2.15亿元，以供应链模型的421.40／2.42亿元只作主体基线检查。两种口径共同说明2025年盈利接近低位，而不是高景气常态。

文莱一期2025年30\%少数股东对应利润约0.73亿元，是主体归属的历史锚。2026Q1集团少数股东损益8.23亿元，占合并净利润约29.2\%，与文莱70\%权益方向一致，但不能据此认定全部少数股东损益来自文莱。半年报若披露非全资子公司财务，是判断文莱利润归属的关键证据。

\begin{exhibitbox}[文莱与集团利润归属基线]
\noindent\begin{tabularx}{\textwidth}{L{3.55cm}R{2.15cm}L{2.7cm}X}
\toprule
\textbf{指标} & \textbf{数值} & \textbf{口径} & \textbf{投资解读} \\
\midrule
恒逸文莱单体2025收入 & 404.74亿元 & 年报主要子公司表 & 项目收入体量可验证 \\
恒逸文莱单体2025净利润 & 2.15亿元 & 年报主要子公司表 & 2025处于低盈利基线 \\
文莱主体基线净利润 & 2.42亿元 & 供应链模型归纳口径 & 只作交叉，不与单体值相加 \\
2025文莱少数股东利润 & 约0.73亿元 & 30\%权益推导 & 高景气必须扣除归属 \\
集团2026Q1总净利润 & 28.18亿元 & 归母19.95+少数8.23 & 总利润不等于归母 \\
集团2026Q1少数股东占比 & 约29.2\% & 8.23／28.18 & 与70\%权益方向一致，非文莱精确拆分 \\
\bottomrule
\end{tabularx}
\end{exhibitbox}
\sourcenote{HY-OFF-2025AR，主要子公司表与少数股东附注；HY-OFF-2026Q1；analysis/supply\_chain\_model.md。主体与单体口径不得相加。}

投资含义是，2026Q1集团利润已经显示少数股东分流不可忽略。任何以文莱税前或合并利润直接推导归母的模型都会高估每股价值；本报告因此在归母层面建模，并把少数股东比例作为下一报告期的核心核验项。

\section{2026H1为何爆发：供给收缩比需求繁荣更重要}

公司在H1预告中称文莱一期满产满销，东南亚成品油及PX、苯盈利处于高位；国内板块亦实现改善。外部行业证据支持“油品利润池上移”的方向：IEA 2026年7月公开摘要显示，成品油裂解价差和炼厂利润处于约四年高位，但其同时预计全球石油需求同比减少约1百万桶／日、炼厂加工量同比减少约2.4百万桶／日。这意味着高利润更多来自炼厂供给收缩与产品紧张，而不是终端需求全面扩张。

公司年报引用的区间显示，2024--2025年PX价差约300--500美元／吨、新加坡柴油裂解价差约15--25美元／桶；2026年公司引用的柴油裂解曾一度超过150美元／桶。该极值来自公司引用第三方口径，未取得可复核逐日Platts序列，不能作为全年平均，更不能直接乘以销量得到净利润。

\begin{exhibitbox}[H1利润爆发的因果链与证据等级]
\small
\noindent\begin{tabularx}{\textwidth}{L{2.8cm}L{2.65cm}L{2.65cm}X}
\toprule
\textbf{因果环节} & \textbf{支持证据} & \textbf{证据等级} & \textbf{仍需验证} \\
\midrule
东南亚产品紧张 & IEA炼厂加工下降、裂解利润高位 & 公共机构摘要，行业方向 & 区域分产品供需与持续时间 \\
文莱高负荷 & 公司预告“满产满销” & 官方、未经审计 & 实际进料、产品产量、检修与库存 \\
柴油／航煤裂解上升 & 公司引用第三方极值、行业方向一致 & 公司转引／行业代理 & 原始逐日序列与公司实现价差 \\
PX／苯盈利上升 & 公司H1原因说明 & 官方定性 & 分产品销量、ASP与成本 \\
税惠与区位 & 先锋企业11年免税、原油近区采购 & 官方制度与公司披露 & 税惠起止、有效税率、实际运费 \\
归母利润释放 & H1归母55--60亿元、扣非54.7--59.7亿元 & 官方预告、未经审计 & 文莱／国内拆分、少数股东与现金流 \\
\bottomrule
\end{tabularx}
\end{exhibitbox}
\sourcenote{HY-OFF-2026H1P；HY-OFF-2025AR；IND-IEA-202607；analysis/industry\_landscape.md，资料截止2026-07-22。公司引用的Platts区间未取得原始序列。}

所以，H1利润反转是真实经营改善，却仍是“周期性经营利润”。扣非接近归母排除了大额非经常性损益主导，但不能证明裂解价差不会均值回归。估值应给予利润真实性信用，同时对持续性打折。

\section{原油、运费、库存和汇率：看价差而不是看油价}

原油约占炼油成本96.2\%，公司披露2025H1／H2原油采购均价约4,074／3,745元／吨。若Brent下降，原料成本有利于后续加工利润，却也可能形成库存损失；若油价上升，成品油价格可能传导，但库存收益和资金占用同时增加。对炼厂而言，绝对油价不是单向利好，综合裂解和库存周转更重要。

外部EIA 2026年7月STEO预计Brent从2026Q2约103美元／桶回落至Q4约70美元／桶、2027年约65美元／桶。该预测是情景而非实现值：若快速回落，文莱可能经历库存收益消退甚至损失，但下游国内聚酯原料压力缓解；若地缘扰动延长，裂解价差可能维持，却伴随更高库存资金和运输风险。

上海航运交易所2026-07-17外部锚显示，中东湾--宁波27万载重吨VLCC运价约61.63美元／吨。文莱原油采购半径更短，公司称供应和运费未受影响；因此该中国长航线运价只能作为压力背景，不能套入文莱实际成本。实际原油贴水、航线、保险、港口费和长约比例仍未披露。

\begin{keyinsight}[文莱一期的“所以呢”]
研究恒逸不是判断Brent单一方向，而是研究“产品裂解减原料、物流、融资和少数股东”后的综合利润。\textbf{行动上，三季度应同时看柴油／航煤裂解、PX价差、文莱负荷、集团少数股东和存货现金流；只看任一单项都不足以提高目标价。}
\end{keyinsight}

\section{税惠：真实优势，但终值证据不足}

文莱一般公司所得税率为18.5\%。公司披露恒逸文莱满足先锋企业条件，可享受11年公司所得税及进口器械、原料免税。税惠能够显著提高高景气利润留存，是项目竞争力的重要组成；但公司没有在本案证据中清晰披露免税具体起止日、2026实际有效税率和到期后安排。

因此，基准模型允许已报告归母利润自然包含当期税惠，却不在终值中额外资本化一项“永久免税价值”。文莱二期的税收优惠批文也不等于一期税惠可无限延长，更不等于二期已经实现相同税率和盈利。

\begin{riskbox}[税惠与归属失效条件]
若半年报显示文莱有效税率高于预期、税惠期限接近结束、少数股东占比显著超过约30\%且无其他解释，或利润未能向母公司形成现金回流，则文莱税后归母能力应下调。\textbf{行动上，先下调H2利润，再审视7.5倍基准PE。}
\end{riskbox}

\section{H2正常化：20亿元是约束，不是悲观口号}

House 2026E归母77.5亿元由H1中值57.5亿元加H2 20.0亿元组成。H2只相当于H1中值约35\%，反映裂解价差回落、可能检修、库存损益、国内需求、未量化转债有效利息和员工计划股份支付等综合缓冲。该处理不是断言文莱利润会崩塌，而是拒绝把异常高价差作为全年平均。

熊／基准／牛H2归母分别为7.5／20.0／37.5亿元。若三季度文莱仍高负荷、区域产品裂解保持高位且现金回流明显，H2可向牛情景移动；若行业裂解迅速回落或检修、运输、税和少数股东折损超预期，则可能向熊情景移动。

\begin{exhibitbox}[H1到H2的正常化桥]
\noindent\begin{tabularx}{\textwidth}{L{2.2cm}R{2.0cm}R{2.0cm}R{2.0cm}X}
\toprule
\textbf{情景} & \textbf{H1中值} & \textbf{H2归母} & \textbf{全年归母} & \textbf{主要假设} \\
\midrule
熊 & 57.5亿元 & 7.5亿元 & 65.0亿元 & 裂解快速回落、现金与费用折损偏大 \\
基准 & 57.5亿元 & 20.0亿元 & 77.5亿元 & 高景气正常化，国内温和修复 \\
牛 & 57.5亿元 & 37.5亿元 & 95.0亿元 & 文莱高价差延续，国内链与现金同步改善 \\
机械年化（拒绝） & 57.5亿元 & 57.5亿元 & 115.0亿元 & 没有周期、检修、归属与费用校准 \\
\bottomrule
\end{tabularx}
\end{exhibitbox}
\sourcenote{HY-OFF-2026H1P；analysis/earnings\_forecast.md；analysis/valuation\_audit.md，预测截止2026-07-22。机械年化仅用于展示错误方法。}

这一章对估值的最终含义是：文莱一期支持2026年高利润，但不支持峰值利润终值化。只有半年报和三季报同时证明负荷、裂解、税后归属与现金流，牛情景才应从20\%概率上调；当前应维持15.7元目标而非向22.5元尾部靠拢。
